% SPDX-License-Identifier: Apache-2.0
\appendix
\section{Statusabgleich und Grenzen der Aussagen}
\label{sec:implementation}
Diese Fassung beschreibt die verbindliche Konzeptkorrektur vom
4. November 2026. Geprüfte Hauptlinie für den jüngsten Konzeptabgleich:
\code{616324e857} (kurzer Git-Präfix). Es wurde ausschließlich Dokumentation
und deren Werkzeugbestand verändert. Die neue globale Anwendungsanerkennung
ist \textbf{noch nicht implementiert}. Ein lesbarer Beleg oder ein erfolgreicher
Dokumentbuild beweist weder Laufzeitkonformität noch Deployment.

\begin{longtable}{@{}P{32mm}P{116mm}@{}}
\toprule
Gegenstand & Nachprüfbare Grenze\\
\midrule\endhead
Abschlussrechnung &
\code{CompletionEvaluator.evaluateAutomaticGraph} stoppt an gemeisterten
Zwischenzielen. In $Q\to B\to A$ mit nur B gemeistert wird A dort
entlassen. Die neue Formel lässt A für die offene Prüfung Q benötigt.
Der historische Profilpfad im selben Service ist ein anderer Algorithmus;
auch seine Existenz beweist keine Umsetzung des neuen Vertrags.\\
Geschützte Ziele & Der bisherige automatische Pfad beginnt bei offenen
Assessments; der neue Vertrag schützt zusätzlich alle ausgewählten
Memorierziele und verlangt eigene fachliche Assessmentabdeckung für jede
gewöhnliche eigene Fachkompetenz. Aggregation und Synthese
werden dadurch nicht zu künstlichen neuen Erwerbsverträgen.\\
Globale Wirkung & Kontextbezogene bisherige Freistellung ist nicht schon
die dauerhaft vertragsgebundene globale Menge $M_{\mathrm{eff}}$.
Folge-Frontier, Gesamt-IST, Vertragsrevision und historische Bedeutung
müssen gemeinsam implementiert und gesondert geprüft werden.\\
Graphgültigkeit & Transitive Minimalität ist verbindlich. Ein lokaler
Validator oder historische Tests ersetzen keine Prüfung der vollständigen
normalisierten Relation und keine fachliche Inhaltsprüfung.\\
Clustererfüllung & Der normative Kern verwendet alle atomaren Nachfahren.
Die gelesenen Backendpfade bevorzugen in bestimmten Fällen rekursiv
\code{isCoreForPrereqs}-Kinder. Das ist nicht dieselbe Aggregation.\\
Memorierung & Der dokumentierte Zieltypvertrag nennt verifizierten Abruf
und keine fälligen Karten. Backend und \code{useSrsMastery} verwenden den
verifizierten Abruf aller erforderlichen Karten. Die Abweichung ist
vor einer entsprechenden Konformitätsbehauptung ausdrücklich aufzulösen.\\
Lerneinheiten & Rekursive Ziel-/Kinderstruktur mit eindeutigem Ersterwerb
und abgeleitetem Abschluss; noch kein behauptetes Runtimefeature.
Technische Ganzheitssperren und direkter Prüfungszugriff unter einem
solchen Profil sind unentschieden. Zeit und Budget bleiben zusätzliche
Planung, kein Synonym für die Einheit.\\
\bottomrule
\end{longtable}

\paragraph{Prüfbare Codeanker.}
\srcref{backend/src/main/java/com/skillpilot/backend/completion/CompletionEvaluator.java}{CompletionEvaluator.java};
\srcref{backend/src/main/java/com/skillpilot/backend/service/LearnerService.java}{LearnerService.java};
\srcref{app/src/hooks/useSrsMastery.ts}{useSrsMastery.ts};
\srcref{app/scripts/validateGraph.ts}{validateGraph.ts};
\srcref{backend/src/main/java/com/skillpilot/backend/domain/Mastery.java}{Mastery.java}.
Die verlinkte Hauptlinie entwickelt sich weiter. Das Build-Manifest bindet
die bei dieser Dokumentfassung tatsächlich gelesenen Dateien per SHA-256.

\section{Prüfstrategie für die Abschlussberechnung}
Das begleitende \code{latex/check-assessment-model.py} ist eine
mathematische Referenzprüfung mit Python-Standardbibliothek. Es liest keine
Lernerdaten und ändert keinen Runtimezustand.

\begin{center}\begin{tabularx}{\linewidth}{@{}P{40mm}X@{}}
\toprule
Prüfbereich & Geprüfte Aussage\\
\midrule
Endliche Modellprüfung & Alle DAGs bis vier Knoten in topologischer
Ordnung mit jeder Mastery-/Schutzmengenbelegung, soweit die lokalen
gewöhnlichen Ziele vollständig assessmentabgedeckt sind:
5\,222 gültige Fälle. Traversierung und statische transitive Hülle
liefern dieselbe implizite Menge.\\
Minimalität & Transitive Reduktion ist minimal und erhält die Hülle.
Redundante Kanten der normalisierten Relation ändern die Anerkennung
nicht; ein Gegenbeispiel prüft die nach Kantenentfernung neu berechnete
Vererbung für verfasste Minimalität.\\
Globale Monotonie & Zusätzliche ausdrückliche Mastery öffnet bei festen
gültigen Verträgen keine effektiv beherrschte Arbeit erneut. Prüfungs-
und Memorierziele werden nie implizit bestanden.\\
Rekursive Einheiten & 128 Belegungen des Mechanikbeispiels prüfen
Abschluss, eigene Anerkennung und explizite Kindpflichten. Ein Aggregat
braucht keine eigene Prüfung; eine reine Synthese braucht keine neuen
gewöhnlichen Ziele.\\
Einheitsgruppierung & 948 Gruppierungsfälle vergleichen den Block-DAG
mit einer vollständigen Blockreihenfolge. Ein Gegenbeispiel zeigt einen
atomaren DAG mit zyklischer Zusammenfassung; keine technische Sperrpolitik
wird daraus abgeleitet.\\
Konkrete Gegenbeispiele & Kein Stopp an Mastery-Zwischenknoten;
lokale Anerkennung trotz späterer Prüfung; dauerhafte Vertragsvereinigung,
künftige Voraussetzungen, Target-Schnitt und externe Unterstützung.\\
Ungültige Eingaben & Zyklen, unbekannte Referenzen und isolierte
gewöhnliche Ziele sind Fehler, keine automatische Beherrschung.\\
\bottomrule
\end{tabularx}\end{center}

Die Diagrammerzeugung prüft die konkret gezeigten Mengen und vergleicht mit
\code{--check} die neu gerenderten PNG-Bytes. Diese Tests prüfen weder
reale Aufgabenqualität noch Benotungszuverlässigkeit, Identitätskontrolle
oder pädagogische Wirkung. Die geforderte unabhängige Inhaltsprüfung
bleibt ein eigener Schritt. Runtimeintegration und deren Sicherheits-,
Nebenläufigkeits- und Oberflächentests sind künftige Implementierungsarbeit.

\section{Begriffe und Symbole}
\begin{longtable}{@{}P{36mm}P{112mm}@{}}
\toprule
Zeichen / Begriff & Bedeutung\\
\midrule\endhead
$G,A,K$ (Kern) & Ziele, strukturelle Atome, Cluster.\\
$C,R,\Reff$ (Kern) & Enthaltensein, direkte und geerbte Voraussetzungen.\\
$R_{\mathrm{at}}$ & Vollständig atomar expandierte effektive Relation;
im Abschlusskapitel kurz $R$.\\
$X^+,X^{-1}$ & Nichtleere transitive Hülle und Umkehrrelation.\\
$\Atoms(g)$ & Atomare Basis; bei einem Atom die einelementige Menge.\\
$m,M$ & Ein gespeicherter binärer Zielzustand und die Menge seiner
ausdrücklichen Erfolge. $m(a)=0$ schließt implizite Beherrschung nicht aus.\\
$M_{\mathrm{eff}},M_A$ & Globale Vereinigung ausdrücklicher und gültig
abgeleiteter impliziter Beherrschung; allgemeine Erfüllungsmenge im Kern.\\
$U,\Own(U),T(U)$ & Lerneinheit, eigene atomare Erwerbsziele und Gesamtumfang einschließlich der Kinder.\\
$\Complete(U)$ & Rekursiv abgeleiteter Abschluss; keine zusätzlich gespeicherte Unit-Mastery.\\
$Q_U,K_U,P_U$ & Eigene Prüfungs-, Memorier- und gemeinsame Schutzmenge.
$K_U$ bezeichnet hier keine strukturellen Cluster.\\
$N_U$ & Vollständige transitive Nachfrage der offenen eigenen geschützten Ziele,
ohne Stopp an gemeisterten Zwischenzielen.\\
$I_U$ & Lokal abgeleitete Anwendung eigener gewöhnlicher Erwerbsziele.
Keine zusätzliche gespeicherte Ergebnisart.\\
Personal Curriculum & Verbindliche persönliche Zielauswahl auf Level 2;
kein Filter, der vorhandene globale Anerkennung zurücknimmt.\\
Fokus & Temporärer Arbeitsbereich auf Personalisierungslevel 3.\\
$T_\pi,B_\pi,\ell_\pi$ & Geplante Erwerbsziele, Eintrittsbasis und
optionale kalenderbezogene Planungsstufen.\\
M7 & Maschineller Curriculum-QA-Reifegrad, kein menschlicher
Wirksamkeitsnachweis.\\
\bottomrule
\end{longtable}

\section{Quellenzuordnung und Abbildungen}
Die Kapitel sind eine redaktionell zusammengeführte deutsche Darstellung
der folgenden Repositoryquellen. Querverweise nennen die jeweilige
fachliche Verantwortung; sie sind keine Behauptung, dass alle optionalen
Konzepte bereits implementiert seien.

\begin{longtable}{@{}P{38mm}P{110mm}@{}}
\toprule
Inhalt & Maßgebliche Quelle\\
\midrule\endhead
Formaler Kern und Filter &
\srcref{docs/concept/skill-graph/graph-definition.md}{Graph Definition},
Abschnitte 1--11.\\
Schichten und Migration &
\srcref{docs/concept/skill-graph/general-goal-system-and-migration.md}{General Goal System and Migration};
\srcref{docs/concept/skill-graph/canonical-gymnasium-rollout.md}{Canonical Gymnasium Rollout}.\\
Sichten und persönliche Auswahl &
\srcref{docs/concept/skill-graph/view-projection-and-goal-placement.md}{View Projection and Goal Placement};
\srcref{docs/concept/levels-of-personalization.md}{Levels of Personalization}.\\
Zieltypen und Evidenz &
\srcref{docs/concept/skill-graph/node-types.md}{Node Types};
\srcref{AGENTS.md}{AGENTS.md, insbesondere Orientierung und Mastery}.\\
Rekursive Lerneinheiten &
\srcref{docs/concept/skill-graph/learning-units.md}{Learning Units};
\srcref{docs/concept/skill-graph/learning-unit-progression.md}{Learning Unit Progression};
\srcref{docs/concept/skill-graph/graph-definition.md}{Graph Definition, Abschnitt 12}.\\
Optionale Zeitplanung &
\srcref{docs/concept/didactic/curriculum-time-axis-and-pacing.md}{Curriculum Time Axis and Pacing}.\\
Impliziter Abschluss &
\srcref{docs/concept/skill-graph/assessment-gated-completion.md}{Assessment-Gated Completion};
\srcref{docs/concept/runtime-workflows/exam-preparation-planning.md}{Exam Preparation Planning}.\\
Herkunft und Materialien &
\srcref{docs/concept/skill-graph/source-and-resource-links.md}{Source and Resource Links};
\srcref{docs/concept/skill-graph/content-integration.md}{Content Integration};
\srcref{docs/concept/skill-graph/human-readable-source-rationales.md}{Human-Readable Source Rationales}.\\
Bilder und Zielbücher &
\srcref{docs/concept/skill-graph/atomic-goal-visualizations.md}{Atomic Goal Visualizations};
\srcref{docs/concept/skill-graph/learning-goal-book-and-evidence-review-pipeline.md}{Learning-Goal Book and Evidence-Review Pipeline}.\\
Validierung und QA &
\srcref{docs/qa-ci/graph-validation-rules.md}{Graph Validation Rules};
\srcref{docs/qa-ci/curriculum-quality-maturity-and-routes.md}{Curriculum Quality Maturity and Routes};
\srcref{docs/qa-ci/semantic-atomicity-review.md}{Semantic Atomicity Review}.\\
Austausch und Rechte &
\srcref{docs/concept/skill-graph/dual-curriculum-package-releases.md}{Dual Curriculum Package Releases};
\srcref{docs/concept/skill-graph/mem-fwu-roundtrip-plan.md}{MEM/FWU Roundtrip Plan};
\srcref{LICENSING.md}{LICENSING.md}.\\
\bottomrule
\end{longtable}

Die sieben didaktischen Graphabbildungen stammen aus
\href{https://github.com/enpasos/skillpilot/tree/main/docs/concept/skill-graph/assets/graph-definition}{SkillPilot: Graph Definition, Abbildungsbestand}
und werden unverändert mit deutschen Erläuterungen verwendet.
Die zusätzlichen Diagramme auf dem Titelblatt sowie zu Zeitplanung,
atomarer Expansion, effektiver Minimalität, transitivem Bedarf, lokaler
Anerkennung, rekursiven Lerneinheiten und Planstand werden deterministisch
mit Matplotlib als PNG erzeugt. Maßgeblich sind die Formeln und Bildlegenden;
die Bilder behaupten keine zusätzlichen Kanten oder Anforderungen.

Technischer Text, Formelspezifikation und Build-Werkzeuge stehen unter
Apache-2.0; die didaktischen Illustrationen unter CC BY 4.0, soweit die
entsprechenden Rechte vorliegen. Lizenztexte und Abgrenzung:
\srcref{LICENSE}{LICENSE},
\srcref{LICENSES/CC-BY-4.0.txt}{LICENSES/CC-BY-4.0.txt} und
\srcref{LICENSING.md}{LICENSING.md}.
Quellenhinweis für die Illustrationen: SkillPilot, enpasos -- Enterprise
Patterns \& Solutions GmbH, Graph Definition, Fassung vom 29. Oktober 2026.
Zusammenführung und deutsche Erläuterungen: diese Dokumentfassung.

\section{Dokument bauen und prüfen}
Die Hauptdatei \code{skillpilot-graph-definition.tex} bindet die Kapitel
aus \code{latex/} und die PNG-Abbildungen aus
\code{assets/graph-definition/} ein. Alle eingebundenen Abbildungen einschließlich Titelbild sind Bitmap-PNGs.

Vom Repositorywurzelverzeichnis aus:
\begin{quote}\scriptsize\ttfamily
python3 docs/concept/skill-graph/latex/generate-completion-diagrams.py --check\par
python3 docs/concept/skill-graph/latex/build-document.py\par
python3 docs/concept/skill-graph/latex/check-assessment-model.py
\end{quote}
Benötigt werden Python 3, Matplotlib und LuaLaTeX mit deutscher Sprachunterstützung
und den Präambelpaketen. Hilfsdateien liegen unter
\code{tmp/skill-graph-document/}, das PDF neben der Hauptdatei.
Der Build löst Verweise in mehreren Läufen auf und bindet Eingabedateien
und PDF durch SHA-256 im Build-Manifest samt Datum und Werkzeugversion.
Über unterschiedliche TeX-Versionen hinweg ist Byteidentität nicht garantiert.

Zum Abschluss gehören ein fehlerfreier LaTeX-Lauf, aufgelöste Verweise,
die Prüfung auf überbreite beziehungsweise überhohe Boxen und eine
visuelle Sichtung des erzeugten PDF. Die Formelprüfung ergänzt diese
Dokumentprüfungen, ersetzt aber weder Runtime-Tests noch unabhängige
fachliche und empirische Validierung.
